% GENERATED FILE. Edit canonical lessons, then run npm run generate:books.
% canonical-source-hash: fnv1a64:fa568a3f9b09f756
% canonical-lessons: TA-C397-makkalatci, TA-C397-pattimanram, TA-C397-orumanataka, TA-C397-tanippatta, TA-C397-pakuttarivu

\chapter{People Power}
\label{ch:ta-people-power}
\hlchaptermodality{397}

\begin{chapteropening}
\textbf{By the end of this chapter:} \emph{I can talk about democracy, debate and reasoned thinking.}

Complete the last lesson of chapter 397: I can talk about democracy, debate and reasoned thinking.
\end{chapteropening}

\section[\texorpdfstring{\ta{மக்களாட்சி} (makkaḷāṭci)}{மக்களாட்சி (makkaḷāṭci)}]{\ta{மக்களாட்சி} (makkaḷāṭci) — democracy}

\label{lesson:TA-C397-makkalatci}



\begin{quote}
Before the new one: say the Tamil for to reflect, then the Tamil for to take into account.
\end{quote}

\subsection*{You'll want to know: \ta{மக்களாட்சி}}
\textbf{\ta{மக்களாட்சி}} — \emph{makkaḷāṭci} — \textquotedblleft{}democracy\textquotedblright{}.

\begin{quote}
\ta{மக்களாட்சியில்} \ta{மக்களே} \ta{தலைவர்களைத்} \ta{தேர்ந்தெடுக்கிறார்கள்}.
\end{quote}

\begin{quote}
\emph{makkaḷāṭciyil makkaḷē talaivarkaḷait tērnteṭukkiṟārkaḷ.}
\end{quote}

\begin{quote}
In a democracy, the people themselves choose their leaders.
\end{quote}

\ta{மக்கள்}, people + \ta{ஆட்சி}, rule: rule by the people.

\begin{cousinweb}
\textbf{\ta{மக்களாட்சி}} (\emph{makkaḷāṭci}) is native: \emph{makkaḷ}, \textquotedblleft{}people\textquotedblright{} (DEDR 4616), and \emph{āṭci}, \textquotedblleft{}rule, government\textquotedblright{}, from \emph{āḷ}, \textquotedblleft{}to rule\textquotedblright{} (DEDR 5157).
\end{cousinweb}

\begin{grammarlens}[title={What you've built}]
One more word for saying what happens if.
\end{grammarlens}

\subsection*{Your turn}
\begin{itemize}
\raggedright
  \item \emph{Say it:} \emph{makkaḷāṭci}
  \item \emph{Say it:} \emph{makkaḷāṭci}, once more
  \item \emph{Say it:} \emph{makkaḷāṭciyil makkaḷē talaivarkaḷait tērnteṭukkiṟārkaḷ.} — In a democracy, the people themselves choose their leaders.
  \item \emph{Recall:} say the Tamil for to reflect, then the Tamil for to take into account, then say \emph{makkaḷāṭci} again
\end{itemize}

\begin{tcolorbox}[breakable,colback=teal!4,colframe=teal!35!black,title={Before you move on}]
What does \ta{மக்களாட்சி} mean? (\textquotedblleft{}Democracy\textquotedblright{}.) Say it once more.
\end{tcolorbox}
\section[\texorpdfstring{\ta{பட்டிமன்றம்} (paṭṭimaṉṟam)}{பட்டிமன்றம் (paṭṭimaṉṟam)}]{\ta{பட்டிமன்றம்} (paṭṭimaṉṟam) — a public debate}

\label{lesson:TA-C397-pattimanram}



\begin{quote}
Before the new one: say the Tamil for to take into account, then the Tamil for democracy.
\end{quote}

\subsection*{You'll want to know: \ta{பட்டிமன்றம்}}
\textbf{\ta{பட்டிமன்றம்}} — \emph{paṭṭimaṉṟam} — \textquotedblleft{}a public debate\textquotedblright{}.

\begin{quote}
\ta{விழாவில்} \ta{ஒரு} \ta{சுவாரசியமான} \ta{பட்டிமன்றம்} \ta{நடந்தது}.
\end{quote}

\begin{quote}
\emph{viḻāvil oru cuvāraciyamāṉa paṭṭimaṉṟam naṭantatu.}
\end{quote}

\begin{quote}
An interesting debate was held at the festival.
\end{quote}

A Tamil stage debate: two teams argue a question before an audience, and a judge gives the verdict.

\begin{cousinweb}
\textbf{\ta{பட்டிமன்றம்}} (\emph{paṭṭimaṉṟam}) is built on native \emph{maṉṟam}, \textquotedblleft{}hall, assembly, court\textquotedblright{} (DEDR 4777); the origin of the first part, \emph{paṭṭi}, is not settled.
\end{cousinweb}

\begin{grammarlens}[title={What you've built}]
One more word for saying what happens if.
\end{grammarlens}

\subsection*{Your turn}
\begin{itemize}
\raggedright
  \item \emph{Say it:} \emph{paṭṭimaṉṟam}
  \item \emph{Say it:} \emph{paṭṭimaṉṟam}, once more
  \item \emph{Say it:} \emph{viḻāvil oru cuvāraciyamāṉa paṭṭimaṉṟam naṭantatu.} — An interesting debate was held at the festival.
  \item \emph{Recall:} say the Tamil for to take into account, then the Tamil for democracy, then say \emph{paṭṭimaṉṟam} again
\end{itemize}

\begin{tcolorbox}[breakable,colback=teal!4,colframe=teal!35!black,title={Before you move on}]
What does \ta{பட்டிமன்றம்} mean? (\textquotedblleft{}A public debate\textquotedblright{}.) Say it once more.
\end{tcolorbox}
\section[\texorpdfstring{\ta{ஒருமனதாக} (orumaṉatāka)}{ஒருமனதாக (orumaṉatāka)}]{\ta{ஒருமனதாக} (orumaṉatāka) — unanimously}

\label{lesson:TA-C397-orumanataka}



\begin{quote}
Before the new one: say the Tamil for democracy, then the Tamil for a public debate.
\end{quote}

\subsection*{You'll want to know: \ta{ஒருமனதாக}}
\textbf{\ta{ஒருமனதாக}} — \emph{orumaṉatāka} — \textquotedblleft{}unanimously\textquotedblright{}.

\begin{quote}
\ta{கூட்டத்தில்} \ta{இந்தத்} \ta{தீர்மானம்} \ta{ஒருமனதாக} \ta{நிறைவேறியது}.
\end{quote}

\begin{quote}
\emph{kūṭṭattil intat tīrmāṉam orumaṉatāka niṟaivēṟiyatu.}
\end{quote}

\begin{quote}
The resolution was passed unanimously at the meeting.
\end{quote}

\ta{ஒரு}, one + \ta{மனம்}, mind + -\ta{ஆக}: with one mind.

\begin{cousinweb}
\textbf{\ta{ஒருமனதாக}} (\emph{orumaṉatāka}) is built from native \emph{oru}, \textquotedblleft{}one\textquotedblright{} (DEDR 990), and \emph{maṉam}, \textquotedblleft{}mind\textquotedblright{}, borrowed from Sanskrit \emph{manas} (CDIAL 9822), with the adverb ending -\ta{ஆக}.
\end{cousinweb}

\begin{grammarlens}[title={What you've built}]
One more word for saying what happens if.
\end{grammarlens}

\subsection*{Your turn}
\begin{itemize}
\raggedright
  \item \emph{Say it:} \emph{orumaṉatāka}
  \item \emph{Say it:} \emph{orumaṉatāka}, once more
  \item \emph{Say it:} \emph{kūṭṭattil intat tīrmāṉam orumaṉatāka niṟaivēṟiyatu.} — The resolution was passed unanimously at the meeting.
  \item \emph{Recall:} say the Tamil for democracy, then the Tamil for a public debate, then say \emph{orumaṉatāka} again
\end{itemize}

\begin{tcolorbox}[breakable,colback=teal!4,colframe=teal!35!black,title={Before you move on}]
What does \ta{ஒருமனதாக} mean? (\textquotedblleft{}Unanimously\textquotedblright{}.) Say it once more.
\end{tcolorbox}
\section[\texorpdfstring{\ta{தனிப்பட்ட} (taṉippaṭṭa)}{தனிப்பட்ட (taṉippaṭṭa)}]{\ta{தனிப்பட்ட} (taṉippaṭṭa) — personal; private}

\label{lesson:TA-C397-tanippatta}



\begin{quote}
Before the new one: say the Tamil for a public debate, then the Tamil for unanimously.
\end{quote}

\subsection*{You'll want to know: \ta{தனிப்பட்ட}}
\textbf{\ta{தனிப்பட்ட}} — \emph{taṉippaṭṭa} — \textquotedblleft{}personal; private\textquotedblright{}.

\begin{quote}
\ta{உங்களுக்குத்} \ta{தனிப்பட்ட} \ta{பிரச்சினை} \ta{இருந்தால்} \ta{என்னிடம்} \ta{சொல்லலாம்}.
\end{quote}

\begin{quote}
\emph{uṅkaḷukkut taṉippaṭṭa piracciṉai iruntāl eṉṉiṭam collalām.}
\end{quote}

\begin{quote}
If you have a personal problem, you can tell me.
\end{quote}

Before a noun only: \ta{தனிப்பட்ட} \ta{முறையில்}, personally, in private.

\begin{cousinweb}
\textbf{\ta{தனிப்பட்ட}} (\emph{taṉippaṭṭa}) is native: \emph{taṉi}, \textquotedblleft{}singleness, solitude\textquotedblright{} (DEDR 3196), with \emph{paṭṭa}, the adjectival past of \emph{paṭu}: \textquotedblleft{}that has become separate\textquotedblright{}.
\end{cousinweb}

\begin{grammarlens}[title={What you've built}]
One more word for saying what happens if.
\end{grammarlens}

\subsection*{Your turn}
\begin{itemize}
\raggedright
  \item \emph{Say it:} \emph{taṉippaṭṭa}
  \item \emph{Say it:} \emph{taṉippaṭṭa}, once more
  \item \emph{Say it:} \emph{uṅkaḷukkut taṉippaṭṭa piracciṉai iruntāl eṉṉiṭam collalām.} — If you have a personal problem, you can tell me.
  \item \emph{Recall:} say the Tamil for a public debate, then the Tamil for unanimously, then say \emph{taṉippaṭṭa} again
\end{itemize}

\begin{tcolorbox}[breakable,colback=teal!4,colframe=teal!35!black,title={Before you move on}]
What does \ta{தனிப்பட்ட} mean? (\textquotedblleft{}Personal; private\textquotedblright{}.) Say it once more.
\end{tcolorbox}
\section[\texorpdfstring{\ta{பகுத்தறிவு} (pakuttaṟivu)}{பகுத்தறிவு (pakuttaṟivu)}]{\ta{பகுத்தறிவு} (pakuttaṟivu) — rational thought; reason}

\label{lesson:TA-C397-pakuttarivu}



\begin{quote}
Before the new one: say the Tamil for unanimously, then the Tamil for personal.
\end{quote}

\subsection*{You'll want to know: \ta{பகுத்தறிவு}}
\textbf{\ta{பகுத்தறிவு}} — \emph{pakuttaṟivu} — \textquotedblleft{}rational thought; reason\textquotedblright{}.

\begin{quote}
\ta{பகுத்தறிவுடன்} \ta{சிந்தித்தால்} \ta{சரியான} \ta{முடிவு} \ta{எடுக்கலாம்}.
\end{quote}

\begin{quote}
\emph{pakuttaṟivuṭaṉ cintittāl cariyāṉa muṭivu eṭukkalām.}
\end{quote}

\begin{quote}
If you think rationally, you can make the right decision.
\end{quote}

\ta{பகுத்து}, having divided + \ta{அறிவு}, knowledge: knowing by taking things apart.

\begin{cousinweb}
\textbf{\ta{பகுத்தறிவு}} (\emph{pakuttaṟivu}) is native: \emph{pakuttu}, \textquotedblleft{}having divided\textquotedblright{}, from \emph{paku} (DEDR 3808), and \emph{aṟivu}, \textquotedblleft{}knowledge\textquotedblright{} (DEDR 314): the knowledge that sorts things out.
\end{cousinweb}

\begin{grammarlens}[title={What you've built}]
One more word for saying what happens if.
\end{grammarlens}

\subsection*{Your turn}
\begin{itemize}
\raggedright
  \item \emph{Say it:} \emph{pakuttaṟivu}
  \item \emph{Say it:} \emph{pakuttaṟivu}, once more
  \item \emph{Say it:} \emph{pakuttaṟivuṭaṉ cintittāl cariyāṉa muṭivu eṭukkalām.} — If you think rationally, you can make the right decision.
  \item \emph{Recall:} say the Tamil for unanimously, then the Tamil for personal, then say \emph{pakuttaṟivu} again
\end{itemize}

\begin{tcolorbox}[breakable,colback=teal!4,colframe=teal!35!black,title={Before you move on}]
What does \ta{பகுத்தறிவு} mean? (\textquotedblleft{}Rational thought; reason\textquotedblright{}.) Say it once more.
\end{tcolorbox}
